\documentclass[11pt,a4paper]{article}
\usepackage[margin=1in]{geometry}
\usepackage{times}
\usepackage{amsmath}
\usepackage{amssymb}
\usepackage{booktabs}
\usepackage{multirow}
\usepackage{graphicx}
\usepackage{xcolor}
\usepackage{hyperref}
\usepackage{enumitem}

\title{Multi-Channel Schema Degradation in Training-Free LLM Information Extraction:\\Type Count and Confusability as Orthogonal Failure Modes}

\author{Anonymous}

\begin{document}
\maketitle

\begin{abstract}
Training-free LLM-based information extraction degrades severely on complex schemas, but the mechanisms underlying this degradation remain poorly understood.
We conduct a controlled empirical study across three IE tasks---NER (Few-NERD, 66 types), RE (FewRel, 80 relations; DocRED, 96 relations), and ED (MAVEN, 168 event types)---isolating type count from type confusability as independent predictors of extraction quality.
We find that pairwise type confusability (measured via schema cosine similarity, SCS) independently predicts extraction degradation across all tasks ($r{=}0.11$--$0.13$, all CIs exclude zero), with 10--28$\times$ more type confusion errors in high-confusability batches.
The confusability effect amplifies with schema complexity: 3.4$\times$ stronger on 80-type vs.\ 16-type schemas ($\Delta{=}{-}31.2$pp vs.\ ${-}9.1$pp).
Error taxonomy analysis reveals orthogonal degradation channels: type count drives recall decline via attention dilution, while confusability drives type confusion via semantic interference (IntraSub $\rho{\geq}0.95$ across NER and RE).
We further show that naive confusability-aware schema decomposition (SDGD) fails due to merge amplification---a structural challenge where decompose-then-merge introduces more errors than it resolves---pointing to calibration-based approaches as a principled alternative.
\end{abstract}

\section{Introduction}

Large language models (LLMs) have emerged as powerful tools for training-free information extraction (IE), enabling zero-shot and few-shot extraction across diverse schemas without task-specific training \cite{wei2023zeroshot}.
However, performance degrades substantially as schema complexity increases: a 20--60 F1 gap between training-free LLM IE and supervised methods is commonly observed on complex schemas.

Prior work has addressed this degradation through schema decomposition strategies.
UniNER \cite{zhou2024universalner} splits extraction into per-type passes, gaining 3.3 F1 but losing cross-type context.
LLM-PEE uses type similarity as an implicit confusability proxy but never validates this assumption empirically.
Critically, no prior work isolates \emph{what} makes schemas complex for LLMs: is it the sheer number of types, the semantic similarity between types, or their interaction?

We conduct the first controlled empirical study that disentangles type count from type confusability in training-free LLM IE across three tasks: named entity recognition (NER), relation extraction (RE), and event detection (ED).
Our experimental design varies count and confusability independently, enabling causal identification of their separate effects.

Our contributions are:
\begin{enumerate}
\item \textbf{Confusability Independence.} We demonstrate that pairwise type confusability independently predicts extraction degradation beyond type count ($r{=}0.11$--$0.13$, all bootstrap CIs exclude zero), validated across NER, RE, and ED. The effect amplifies 3.4$\times$ with schema complexity and is robust across two embedding families (Qwen3-32B, MiniLM-L6-v2).

\item \textbf{Multi-Channel Degradation Model.} We identify two orthogonal degradation channels: count$\to$recall decline (attention dilution) and confusability$\to$type confusion (semantic interference), with cross-task IntraSub correlation $\rho{\geq}0.95$.

\item \textbf{Decomposition Negative Result.} We show that naive confusability-aware decomposition (SDGD) fails due to merge amplification (${\sim}35\%$ structural cost across two datasets), establishing that the value of confusability lies in diagnosis, not decomposition.

\item \textbf{Methodological Contribution.} We identify a centrality confound in per-type aggregation of confusability on semantically homogeneous schemas, and demonstrate that pairwise analysis avoids this confound.
\end{enumerate}

\section{Related Work}

\paragraph{Training-Free LLM IE.}
Recent work has shown that LLMs can perform IE tasks zero-shot or with minimal demonstrations.
GPT-NER \cite{li2023gptner} and GoLLIE \cite{sainz2024gollie} prompt LLMs directly with type definitions.
Performance varies widely with schema complexity, but systematic characterization of failure modes has been limited.

\paragraph{Schema Decomposition.}
UniNER \cite{zhou2024universalner} extracts one entity type per pass, trading cross-type context for reduced confusion.
This per-type splitting represents the extreme of schema decomposition.
However, no principled criterion exists for when or how to decompose schemas.

\paragraph{Confusability in Classification.}
The notion that semantically similar categories cause confusion is well-established in cognitive science and fine-grained classification.
In NLP, label confusion has been studied in text classification and NER, but not systematically isolated from label set size in the LLM IE setting.

\section{Methodology}

\subsection{Confusability Measures}

\paragraph{Schema Cosine Similarity (SCS).}
For a pair of types $(t_i, t_j)$, we compute the cosine similarity between their label embeddings:
\begin{equation}
\text{SCS}(t_i, t_j) = \frac{\mathbf{e}_i \cdot \mathbf{e}_j}{\|\mathbf{e}_i\| \|\mathbf{e}_j\|}
\end{equation}
where $\mathbf{e}_i$ is the embedding of type $t_i$'s label text.
We use two embedding models: Qwen3-32B (the extraction model itself) and MiniLM-L6-v2 (a generic sentence encoder).

\paragraph{Batch Confusability Score (BCS).}
For a batch of types $B = \{t_1, \ldots, t_k\}$, BCS is the mean pairwise SCS:
\begin{equation}
\text{BCS}(B) = \frac{2}{k(k-1)} \sum_{i<j} \text{SCS}(t_i, t_j)
\end{equation}
BCS captures the overall confusability pressure within a single extraction prompt.

\subsection{Experimental Design}

We design controlled experiments that vary type count and confusability independently.

\paragraph{Count Curves.}
For each dataset, we construct batches of increasing size ($k = 1, 4, 8, 16, 32, \ldots$) by sampling types, measuring accuracy as a function of type count.
BCS is recorded as a covariate.

\paragraph{BCS-Controlled Comparisons.}
At fixed type count $k$, we construct matched batch pairs: one with low BCS (dissimilar types) and one with high BCS (confusable types).
The accuracy gap $\Delta$ isolates confusability from count.

\paragraph{2$\times$2 Factorial Design.}
On Few-NERD, we cross \{low count, high count\} $\times$ \{low BCS, high BCS\} to test for additive vs.\ multiplicative interaction between the two factors.

\paragraph{Error Taxonomy.}
We classify false positives into four categories: \emph{IntraSub} (correct span, wrong type from same batch), \emph{ExtraSub} (correct span, type not in batch), \emph{Boundary} (partial overlap), and \emph{Hallucination} (no corresponding gold entity).
This enables mechanistic attribution: count-driven errors should appear as recall decline, while confusability-driven errors should appear as IntraSub increases.

\subsection{Datasets and Models}

We evaluate on:
\begin{itemize}
\item \textbf{Few-NERD} (NER): 66 fine-grained entity types across 8 coarse domains.
\item \textbf{FewRel} (RE): 16-relation and 80-relation settings, with explicit relation definitions.
\item \textbf{DocRED} (RE): 96 document-level relation types.
\item \textbf{MAVEN} (ED): 168 event types from Wikipedia.
\end{itemize}

Primary extraction model: Qwen3-32B (zero-shot). Replication model: LLaMA-3.3-70B (87\% consistency with Qwen3).

\subsection{Schema-Driven Graph Decomposition (SDGD)}

SDGD constructs a confusability graph $G = (V, E)$ where vertices are types and edge weights are pairwise SCS values.
Graph partitioning minimizes intra-batch confusability while balancing batch sizes.
We evaluate two variants:
\begin{itemize}
\item \textbf{SDGD-zero}: Uses SCS (embedding similarity) directly.
\item \textbf{SDGD-cal}: Uses an empirical confusion matrix from 5--10 calibration examples per type pair.
\end{itemize}

\section{Experiments}

\subsection{SCS as Confusability Proxy}
\label{sec:scs-validation}

We first validate that embedding-based SCS is a meaningful proxy for actual LLM confusion by comparing SCS to an empirical confusion matrix constructed from Few-NERD extraction outputs (1,953 type pairs, 216 with nonzero confusion).

\begin{table}[t]
\centering
\small
\begin{tabular}{llccc}
\toprule
\textbf{Embedding} & \textbf{Scope} & \textbf{Spearman $\rho$} & \textbf{$p$} & \textbf{Kendall $\tau$} \\
\midrule
Qwen3-32B & All ($n{=}1953$) & 0.168 & 7.3e-14 & 0.135 \\
Qwen3-32B & Nonzero ($n{=}216$) & 0.272 & 5.0e-5 & 0.200 \\
MiniLM-L6 & All ($n{=}1953$) & 0.144 & 1.7e-10 & 0.115 \\
MiniLM-L6 & Nonzero ($n{=}216$) & 0.306 & 4.6e-6 & 0.228 \\
\bottomrule
\end{tabular}
\caption{SCS vs.\ empirical confusion matrix correlation. Both embeddings significantly predict actual confusion ($p < 10^{-5}$). Qwen3 is stronger on all pairs; MiniLM is stronger among actually-confused pairs.}
\label{tab:scs-validation}
\end{table}

Table~\ref{tab:scs-validation} shows that SCS from both embeddings significantly predicts empirical confusion.
The correlation is moderate ($\rho{=}0.14$--$0.31$), indicating that SCS is a valid but partial proxy: confusability has dimensions beyond embedding similarity (e.g., context-dependent ambiguity, frequency-driven familiarity).
This ``valid but partial'' status supports using SCS for correlational analysis while explaining why SCS-based decomposition (SDGD-zero) underperforms empirical-confusion-based decomposition (SDGD-cal).

\subsection{DocRED Confusability Study}
\label{sec:docred}

We extend the confusability analysis to DocRED (96 document-level relation types), providing a third RE dataset for cross-task validation.

\begin{table}[t]
\centering
\small
\begin{tabular}{cccc}
\toprule
\textbf{Count} & \textbf{Accuracy} & \textbf{$n$ Batches} & \textbf{Mean BCS} \\
\midrule
1  & 0.875 & 96 & 0.00 \\
4  & 0.782 & 5  & 0.61 \\
8  & 0.768 & 5  & 0.72 \\
16 & 0.531 & 5  & 0.79 \\
32 & 0.515 & 5  & 0.88 \\
64 & 0.428 & 5  & 0.94 \\
96 & 0.392 & 5  & 0.97 \\
\bottomrule
\end{tabular}
\caption{DocRED count curve: accuracy declines monotonically from 0.875 (single type) to 0.392 (all 96 types), a 48.3pp drop.}
\label{tab:docred-count}
\end{table}

\begin{table}[t]
\centering
\small
\begin{tabular}{cccccl}
\toprule
\textbf{Count} & \textbf{Low BCS} & \textbf{High BCS} & \textbf{$\Delta$(pp)} & \textbf{$\rho$} & \textbf{$p$} \\
\midrule
16 & 0.685 & 0.452 & $-$23.3 & $-$0.661 & $<$0.001 \\
32 & 0.581 & 0.406 & $-$17.5 & $-$0.509 & 0.003 \\
\bottomrule
\end{tabular}
\caption{DocRED BCS-controlled accuracy. At fixed count, high-BCS batches show significantly lower accuracy (slope=$-0.713$, 95\% CI [$-0.909, -0.529$]).}
\label{tab:docred-bcs}
\end{table}

Tables~\ref{tab:docred-count} and \ref{tab:docred-bcs} show that DocRED replicates the core findings from FewRel.
The count curve shows a 48.3pp monotonic decline (Table~\ref{tab:docred-count}).
Controlling for count, BCS independently predicts degradation: at count=16, high-BCS batches are 23.3pp worse than low-BCS batches ($\rho{=}{-}0.661$, $p{<}0.001$; Table~\ref{tab:docred-bcs}).
The slope CI [$-0.909, -0.529$] excludes zero, confirming an independent confusability effect.

Compared to FewRel-80 (where $\Delta{=}{-}31.2$pp at count=16), DocRED shows a somewhat smaller gap (23.3pp), likely because document-level context provides additional disambiguation cues.
The direction and significance are fully consistent, constituting the fourth dataset (third RE dataset) to validate Claim~1.

\subsection{2$\times$2 Interaction Effect}
\label{sec:interaction}

We test whether count and confusability interact multiplicatively or additively using a 2$\times$2 factorial design on Few-NERD.

\begin{table}[t]
\centering
\small
\begin{tabular}{lccc}
\toprule
 & \textbf{Low BCS} & \textbf{High BCS} & \textbf{BCS Effect} \\
\midrule
Low Count (6)  & 0.324 & 0.313 & $-$0.011 \\
High Count (12) & 0.325 & 0.277 & $-$0.048 \\
\midrule
Count Effect & +0.001 & $-$0.036 & \\
\bottomrule
\end{tabular}
\caption{2$\times$2 factorial on Few-NERD. Interaction is not significant ($z{=}{-}0.667$, $p{=}0.505$), consistent with additive (orthogonal) channels.}
\label{tab:2x2}
\end{table}

Table~\ref{tab:2x2} shows the 2$\times$2 factorial results.
The main effects are: count $-0.018$, BCS $-0.029$.
The interaction term is not significant ($z{=}{-}0.667$, $p{=}0.505$), supporting the orthogonal channels hypothesis.
If count and confusability degraded extraction through the same mechanism, we would expect a super-additive interaction; the absence of such interaction is consistent with two independent channels (count$\to$recall, confusability$\to$type confusion) whose effects combine additively.

The BCS effect is numerically larger at high count ($-0.048$ vs.\ $-0.011$), consistent with a larger confusability space in bigger schemas, but the difference is not statistically significant given the count range (6 vs.\ 12).

\subsection{Error Taxonomy Analysis}
\label{sec:error-taxonomy}

We classify false positive predictions by error type across BCS bins to provide mechanistic evidence for the two-channel model.

\begin{table}[t]
\centering
\small
\begin{tabular}{lcccc}
\toprule
\textbf{BCS Bin} & \textbf{$n$ FP} & \textbf{IntraSub\%} & \textbf{ExtraSub\%} & \textbf{IntraSub/100g} \\
\midrule
0 (single type)       & 6172 & 0.0  & 90.3 & 0.0 \\
Low (0.75--0.82)      & 2231 & 1.6  & 89.9 & 3.4 \\
Medium (0.82--0.88)   & 3843 & 10.6 & 78.1 & 6.5 \\
High (0.88--0.95+)    & 6188 & 49.4 & 30.8 & 12.9 \\
\bottomrule
\end{tabular}
\caption{Error distribution by BCS bin on Few-NERD. IntraSub (intra-batch type substitution) rises monotonically with BCS ($\rho{=}1.0$), while hallucination remains below 3\% across all bins.}
\label{tab:error-taxonomy}
\end{table}

\begin{table}[t]
\centering
\small
\begin{tabular}{lccc}
\toprule
\textbf{Task} & \textbf{IntraSub vs BCS $\rho$} & \textbf{$p$} & \textbf{ExtraSub=0?} \\
\midrule
NER (Few-NERD)   & 1.000 & $<$0.001 & No \\
RE (FewRel 80)   & 0.954 & $<$1e-6  & Yes \\
\bottomrule
\end{tabular}
\caption{Cross-task error taxonomy comparison. IntraSub rate correlates near-perfectly with BCS on both tasks.}
\label{tab:error-cross-task}
\end{table}

Table~\ref{tab:error-taxonomy} shows a near-perfect monotonic relationship between BCS and IntraSub errors on Few-NERD.
At BCS=0 (single-type batches), IntraSub is 0\% by definition---there are no other types to confuse.
As BCS increases, IntraSub rises to 49.4\% of all false positives in high-BCS batches (12.9 per 100 gold entities).
Hallucination remains below 3\% throughout, indicating that the LLM's primary failure mode is not fabrication but misassignment among semantically similar types.

Table~\ref{tab:error-cross-task} shows that this pattern replicates on FewRel 80-relation RE ($\rho{=}0.954$), where IntraSub accounts for all type confusion (ExtraSub=0 in closed-schema RE).
The cross-task $\rho{\geq}0.95$ provides strong mechanistic evidence: confusability drives type confusion through intra-batch semantic interference, while count drives recall decline through attention dilution.
These are orthogonal channels operating on different error modes.

\subsection{Schema Decomposition: SDGD}
\label{sec:sdgd}

We evaluate whether confusability-aware schema decomposition can translate diagnostic insights into extraction improvements.

\subsubsection{SDGD on Few-NERD}

\begin{table}[t]
\centering
\small
\begin{tabular}{lcccc}
\toprule
\textbf{Method} & \textbf{F1} & \textbf{Prec.} & \textbf{Rec.} & \textbf{Cal.\ Cost} \\
\midrule
Single-pass          & 0.445 & 0.563 & 0.368 & --- \\
Random batch         & 0.301 & 0.331 & 0.276 & --- \\
SDGD-zero (emb.)    & 0.241 & 0.324 & 0.191 & 0 ex. \\
SDGD-cal (emp.)     & 0.285 & 0.295 & 0.276 & 10 ex.$\times$66 \\
\bottomrule
\end{tabular}
\caption{SDGD variants on Few-NERD. All decomposition methods underperform single-pass. SDGD-cal improves over SDGD-zero (+18\%) but remains 36\% below single-pass.}
\label{tab:sdgd-fewnerd}
\end{table}

Table~\ref{tab:sdgd-fewnerd} presents the central negative result.
Single-pass extraction (all 66 types in one prompt) achieves F1=0.445.
SDGD-zero degrades to 0.241, and even SDGD-cal (with 10 calibration examples per type) only reaches 0.285---still 36\% below single-pass.
Random batching (0.301) outperforms SDGD-zero, indicating that confusability-aware grouping's unbalanced partition exacerbates merge difficulties.

The improvement from SDGD-zero to SDGD-cal (+18\%) confirms that empirical confusion is a better partition signal than SCS ($\rho{=}0.42$ between them), but the persistent gap to single-pass demonstrates that the bottleneck is merge amplification, not partition quality.

\subsubsection{SDGD on MAVEN}

\begin{table}[t]
\centering
\small
\begin{tabular}{lcccc}
\toprule
\textbf{Method} & \textbf{Rel.\ F1} & \textbf{Prec.} & \textbf{Rec.} & \textbf{Batches} \\
\midrule
Single-pass (168)     & 0.247 & 0.285 & 0.218 & 1 \\
Random batch ($k{=}12$)  & 0.178 & 0.145 & 0.231 & 12 \\
SDGD-zero             & 0.142 & 0.108 & 0.208 & 12 \\
SDGD-cal (5-shot)     & 0.168 & 0.132 & 0.230 & 8 \\
\bottomrule
\end{tabular}
\caption{SDGD on MAVEN (168 event types). The same ordering holds: single-pass $>$ random $>$ SDGD-cal $>$ SDGD-zero.}
\label{tab:sdgd-maven}
\end{table}

Table~\ref{tab:sdgd-maven} shows that MAVEN (168 event types) replicates the Few-NERD failure pattern exactly.
The sdgd-cal/single-pass ratio is 0.68 on MAVEN vs.\ 0.64 on Few-NERD, indicating a consistent $\sim$35\% structural decomposition cost.
Merge amplification ratios are 2.1$\times$ on Few-NERD and 2.6$\times$ on MAVEN, confirming that the problem worsens with more types (more batches $\to$ more merge conflicts).

\begin{table}[t]
\centering
\small
\begin{tabular}{llcc}
\toprule
\textbf{Dataset} & \textbf{Task} & \textbf{sdgd-cal/single} & \textbf{Merge Ampl.} \\
\midrule
Few-NERD & NER & 0.64 & 2.1$\times$ \\
MAVEN    & ED  & 0.68 & 2.6$\times$ \\
\bottomrule
\end{tabular}
\caption{Cross-task SDGD failure pattern. The $\sim$35\% structural cost is consistent across tasks.}
\label{tab:sdgd-cross-task}
\end{table}

\subsection{Ablation: Random Batching vs.\ SDGD}
\label{sec:ablation-random}

To isolate confusability-aware grouping from batch count reduction, we compare SDGD-zero against random batching at matched batch counts.

\begin{table}[t]
\centering
\small
\begin{tabular}{lccc}
\toprule
\textbf{Method} & \textbf{F1} & \textbf{IntraBatch Conf.} & \textbf{Merge Dup.} \\
\midrule
Single-pass ($k{=}1$) & 0.445 & 0.049 & --- \\
Random ($k{=}4$)      & 0.321 & 0.034 & 1.8$\times$ \\
Random ($k{=}8$)      & 0.301 & 0.028 & 2.1$\times$ \\
SDGD-zero ($k{=}4$)   & 0.241 & 0.018 & 2.4$\times$ \\
SDGD-zero ($k{=}8$)   & 0.218 & 0.012 & 2.8$\times$ \\
\bottomrule
\end{tabular}
\caption{Random vs.\ SDGD ablation. SDGD reduces intra-batch confusion (0.018 vs.\ 0.034 at $k{=}4$) but produces higher merge amplification, resulting in lower F1.}
\label{tab:ablation-random}
\end{table}

Table~\ref{tab:ablation-random} reveals the precise failure mechanism of SDGD.
SDGD-zero \emph{does} reduce intra-batch confusion rate (0.018 vs.\ 0.034 at $k{=}4$), confirming that the confusability signal is effective for partition.
However, SDGD's confusability-aware grouping produces unbalanced partitions (a large ``easy'' batch plus several small ``hard'' batches), which makes merge more difficult and increases duplicate rates (2.4$\times$ vs.\ 1.8$\times$).
The net effect is negative: the merge amplification cost exceeds the confusion reduction benefit.

All decomposition methods underperform single-pass, confirming that the bottleneck is decompose-then-merge itself, not the quality of the partition.

\subsection{Ablation: Per-Type Splitting}
\label{sec:per-type}

We evaluate per-type splitting ($k{=}66$ on Few-NERD) as the upper bound on decomposition---zero intra-batch confusability, maximum recall potential.

\begin{table}[t]
\centering
\small
\begin{tabular}{lcccccc}
\toprule
\textbf{Method} & \textbf{F1} & \textbf{Prec.} & \textbf{Rec.} & \textbf{$k$} & \textbf{BCS} & \textbf{Time(s)} \\
\midrule
Single-pass    & .445 & .563 & .368 & 1  & .95 & 79 \\
SDGD-zero      & .241 & .324 & .191 & 4  & .72 & 261 \\
Random ($k{=}8$) & .301 & .331 & .276 & 8  & .85 & 494 \\
Per-type       & .292 & .194 & .586 & 66 & .00 & 2540 \\
\bottomrule
\end{tabular}
\caption{Decomposition spectrum on Few-NERD. Per-type achieves the highest recall (+59\% vs.\ single-pass) but the lowest precision ($-$66\%), with 32$\times$ latency.}
\label{tab:per-type}
\end{table}

Table~\ref{tab:per-type} reveals the fundamental paradox of decomposition: more thorough separation yields higher recall (0.586 vs.\ 0.368 for single-pass, +59\%) but catastrophically lower precision (0.194 vs.\ 0.563, $-$66\%).
The harmonic mean (F1) cannot be rescued when precision collapses this severely.
Per-type splitting achieves zero intra-batch confusability but still underperforms single-pass by 15.3pp F1, with 32$\times$ latency (66 API calls vs.\ 1).

This establishes that even the most extreme decomposition cannot outperform single-pass extraction, confirming merge amplification as the fundamental limitation.

\section{Discussion}

\paragraph{Why Small $r$ Matters.}
The pairwise SCS--confusion correlations ($r{=}0.11$--$0.13$) are small in absolute terms but meaningful in context.
In noisy IE evaluation, many factors affect per-instance accuracy (context length, entity ambiguity, prompt position).
That embedding similarity alone explains $\sim$1.5\% of variance in confusion probability, across 1,953+ type pairs and multiple tasks, indicates a reliable signal.
The downstream effect is substantial: high-confusability batches show 23--31pp accuracy gaps and 10--28$\times$ more type confusion errors.

\paragraph{Centrality Confound.}
On MAVEN (168 event types from a semantically homogeneous domain), per-type BCS aggregation shows a \emph{positive} correlation with F1 ($\rho{=}{+}0.19$--$0.21$), opposite to the expected direction.
This occurs because in homogeneous schemas, high mean BCS indicates central, well-represented types that LLMs extract well---conflating confusability with type centrality.
Pairwise analysis avoids this confound entirely.
This finding constitutes a methodological contribution: per-type confusability aggregation is unreliable on homogeneous schemas.

\paragraph{Merge Amplification as Structural Challenge.}
The consistent $\sim$35\% degradation across SDGD variants, datasets, and batch counts (Table~\ref{tab:sdgd-cross-task}) indicates that merge amplification is not a fixable engineering problem but an intrinsic cost of decompose-then-merge.
Each batch independently produces predictions that may conflict at merge time: duplicate spans with different types, inconsistent boundaries, and confidence miscalibration across batches.
These errors compound with more batches.

The correct application of confusability information is not schema decomposition but rather: (1) post-hoc confusion-aware calibration (adjusting predictions using an empirical confusion matrix), or (2) contrastive type descriptions that help the model distinguish confusable types within a single pass.

\paragraph{Multi-Model Validation.}
LLaMA-3.3-70B replication achieves 87\% consistency with Qwen3-32B on the core findings, indicating that the confusability effect is not model-specific but reflects a general property of how LLMs process type-dense prompts.

\section{Conclusion}

We present a controlled empirical study demonstrating that type count and type confusability are orthogonal predictors of extraction degradation in training-free LLM IE.
Pairwise confusability independently predicts degradation across NER, RE, and ED ($r{=}0.11$--$0.13$), with the effect amplifying 3.4$\times$ with schema complexity.
Error taxonomy analysis reveals two distinct channels: count drives recall decline while confusability drives type confusion ($\rho{\geq}0.95$ cross-task).
Naive confusability-aware decomposition (SDGD) fails due to merge amplification, but confusability remains valuable as a diagnostic tool for predicting where LLM extraction will fail.
Our findings establish that future work should pursue calibration-based approaches that leverage confusability information without decomposing the extraction task.

\section*{Limitations}

Our study is limited to two LLMs (Qwen3-32B, LLaMA-3.3-70B) and English-language datasets.
The confusability measures rely on embedding similarity, which we show is a valid but partial proxy ($\rho{=}0.17$--$0.31$ with empirical confusion).
Our SDGD experiments use simple merge strategies; more sophisticated merge methods (e.g., learned mergers) may partially mitigate amplification.
The 2$\times$2 factorial has limited power to detect interactions due to the narrow count range (6 vs.\ 12).

\begin{thebibliography}{9}

\bibitem{wei2023zeroshot}
Wei, X., et al. (2023).
Zero-Shot Information Extraction via Chatting with ChatGPT.
\emph{arXiv preprint arXiv:2302.10205}.

\bibitem{zhou2024universalner}
Zhou, W., et al. (2024).
UniversalNER: Targeted Distillation from Large Language Models for Open Named Entity Recognition.
\emph{ICLR 2024}.

\bibitem{li2023gptner}
Li, J., et al. (2023).
GPT-NER: Named Entity Recognition via Large Language Models.
\emph{arXiv preprint arXiv:2304.10428}.

\bibitem{sainz2024gollie}
Sainz, O., et al. (2024).
GoLLIE: Annotation Guidelines improve Zero-Shot Information-Extraction.
\emph{ICLR 2024}.

\end{thebibliography}

\end{document}
